\scanpage{49}
% L1, heading
Conditions à envisager

% L2
(i) $\underline C\xrightarrow{i_*}\underline A$ transforme injectifs en injectifs

% L2–7, brace beside (i)–(iii)
% slanted
\margin{\unsure{équivalentes|équivalents} $\left\{\vphantom{\rule{0pt}{4em}}\right.$}

% L3–5
(ii) Pour tout objet injectif $x$ de $\underline A$, le plus grand sous-objet $x'$ de $x$ qui est $\in\mathrm{Ob}\,\underline C$ est encore injectif \uline{dans} $A$.

% L6–8
(iii) Pour tout objet de $\underline C$, son enveloppe injective dans $\underline A$ est dans $\underline C$ [i.e. $\underline C$ stable par enveloppes injectives] -- donc \unsure{néc.|aussi|a. c.} injective dans $\underline C$.

% L8–14, two arrows below the brace, ending at (iv) and (v)
\diagram{\begin{tikzcd}[ampersand replacement=\&,row sep=1.2em,column sep=1em]
{} \arrow[d] \arrow[dd,bend right=22,Rightarrow] \& \\ {} \& \\ {} \&
\end{tikzcd}}

% L9–13
(iv) \deleted{Pour tout $x$} L'hom naturel
\[
R^0(\mathrm{id}/\unsure{i_*|[{}_*}i^!)\longrightarrow ST
\]
est un isomorphisme, i.e. \deleted{pour} $x$ objet injectif de $\underline A$, \deleted{$\unsure{x/i|x/i^{!}|\text{si}}$} si $x'$ \inserted{$=i_*i^{\prime !}(x')$} est le plus grand sous-objet $\in\mathrm{Ob}\,\underline C$, alors $x/x'$ est $C''$-fermé.

% L12–14, margin crossing the longer arrow
% slanted
\margin{cf Gabriel, cor. 3}

% L14
(v) $T:\underline A\unsure{\mathord{-}|\to}\underline A/\underline C$ transforme injectifs en injectifs.

% L15–17
$(\mathrm{i})\Rightarrow(\mathrm{ii})$ Car \deleted{$x'=$} pour \uline{tt} $x$, $x'=i_*i^!(x)$, or $i^!(x)$ [étant adjoint de $i_*$ qui est exact] transforme injectifs en injectifs

% L18–22
$(\mathrm{ii})\Rightarrow(\mathrm{i})$ Car si $x\in\underline C$, soit $\bar x$ une enveloppe injective de $x$ dans $\underline A$, \deleted{et} considérons \deleted{$\unsure{i_*|i}$}$\bigl(i^!(\bar x)\bigr)$, c'est un objet \uline{injectif} dans $\underline C$ contenant $x$, et une extension essentielle du dit [car elle l'est dans $C$] donc c'est une enveloppe injective de $x$ \uline{dans} $\underline C$.

% L23–24
Si $\bar x$ est injectif, il est \unsure{en} égal à $i^!(\bar x)$. Si (ii) est vérifié, $i_*i^!(\bar x)$ est injectif, i.e. $i_*(x)$ injectif.

% archive sheet number: 48
